\section{Warming at the critical time scale}\label{sec:warming}

The first step toward a quench limit is to control the finite law after
any positive multiple of \(\Tn=n^{2/3}\), uniformly in its initial
distribution.  The edge estimates of Proposition~\ref{prop:edge-estimates} give the
required control at the smallest observation scales.  We turn them into
dissipation estimates by measuring how much relative entropy survives
between two observations.  Larger entropy levels are treated by the
observation path to a product measure from \cite[Section~19]{CM}.

For a probability \(\rho\) on the cube, write
\[
 w_s^\rho=\frac{\dd(\rho P_{s\Tn})}{\dd\mu},\qquad s\ge0.
\]
Thus \(w_s=w_s^{\nu_n}\) for the uniform initial law used elsewhere in
the paper.  Write
\[
 p_t(\sigma,\tau)=\frac{P_t(\sigma,\tau)}{\mu(\tau)}
\]
for the heat kernel relative to \(\mu\).  We use
\(\log_+z=\max\{\log z,0\}\) for \(z>0\), with
\(z(\log_+z)^m=0\) at \(z=0\).
For a probability \(\pi\) and a nonnegative integrable function \(z\), write
\[
 \Ent_\pi(z)=\pi[z\log z]-\pi[z]\log\pi[z],
\]
with \(0\log0=0\) and \(\Ent_\pi(0)=0\).  For a probability
\(\nu\ll\pi\), let
\(D(\nu\Vert\pi)=\Ent_\pi(\dd\nu/\dd\pi)\), and set this relative
entropy to \(+\infty\) otherwise.

\begin{theorem}[Warming from arbitrary initial laws]\label{thm:warming}
Suppose that the couplings are Gaussian.  Fix \(\vartheta>0\), an integer
\(m\ge1\), and \(0<\xi<1/100\).  There are measurable disorder events
\(\mathcal G_n\), a constant \(C<\infty\), and deterministic
\(\varepsilon_n\downarrow0\) such that
\(\liminf_{n\to\infty}\Prob(\mathcal G_n)\ge1-\vartheta\), and the
following inequalities hold on \(\mathcal G_n\), for all sufficiently
large \(n\).  The constant is independent of the initial law \(\rho\),
which may be chosen after the disorder.
\begin{align}
 \Ent_\mu(w_s^\rho)&\le C(1+s^{-1.54}),
       &&s>0,\quad\hbox{for every probability }\rho ,
       \label{warm:entropy}\\
 \|P_{s\Tn}\|_{L^1(\mu)\to L^\infty(\mu)}
       &=\max_{\sigma,\tau}p_{s\Tn}(\sigma,\tau)
       \le C\exp(C s^{-1.54}),
       &&s\ge\varepsilon_n ,
       \label{warm:kernel}\\
 \left(1+\mu\!\left[
       w_s^\rho(\log_+w_s^\rho)^m\right]\right)^{1/m}
       &\le C\left(1+s^{-1/(2/3-\xi)}\right),
       &&s>0,\quad\hbox{for every probability }\rho .
       \label{warm:log-moment}
\end{align}
The events and constants may depend on the displayed fixed parameters.
In particular the order \(m\) is fixed before \(n\) tends to infinity.
\end{theorem}

The distinction between \eqref{warm:entropy} and
\eqref{warm:kernel} is useful.  Entropy is controlled at every positive
time; the stronger kernel estimate begins at a deterministic time tending
to zero in the rescaled units.  The log-moment estimate supplies a
different form of control at the very short times used later in the
replacement argument.

\subsection{A restricted dissipation profile}\label{warm:profile-section}

Recall \(\cI(z)=\cE(z,\log z)\), with the extended value at zeros.  On a
two-point edge, the inequality
\[
 (a^2-b^2)(\log a^2-\log b^2)\ge4(a-b)^2,\qquad a,b>0,
\]
follows by integrating \(1/u\) between \(a\) and \(b\) and applying
Cauchy--Schwarz.  Summing it in the reversible edge representation gives
\begin{equation}
 \cI(f^2)\ge4\cE(f).                                      \label{warm:modified-classical}
\end{equation}
The inequality extends to nonnegative \(f\) by decreasing positive
regularization.

The next profile concerns a bounded function whose square tilt retains
a definite fraction of the logarithm of its inverse mass.  All its
inequalities hold simultaneously in the function.

\begin{proposition}[Restricted dissipation profiles]\label{warm:profiles}
Fix \(C_{\mathrm{lev}}<\infty\) and \(0<\xi<1/100\).  With tight constants, there
are a threshold \(K_0<\infty\) and a number \(\delta>0\) such that every
\(0\le f\le1\) satisfying
\begin{equation}
 N=\mu[f^2]=e^{-L},\qquad
 K_0\le L\le C_{\mathrm{lev}} n,\qquad
 D\!\left(\frac{f^2\mu}{N}\middle\Vert\mu\right)\ge\frac L2
                                                               \label{warm:restricted-class}
\end{equation}
obeys
\begin{align}
 \Tn\,\frac{\cI(f^2)}{N}&\ge c_\xi L^{5/3-\xi},
                                                               \label{warm:modified-profile}\\
 \Tn\,\frac{\cE(f)}{N}&\ge c
 \begin{cases}
 L^{1.662},&L\le n^\delta,\\
 L^{1.162},&L>n^\delta .
 \end{cases}                                                   \label{warm:classical-profile}
\end{align}
Here a threshold, a positive exponent, and constants are tight in the
following precise sense.  For each \(\vartheta>0\) they can be fixed,
independently of \(f\), on disorder events of limiting lower probability
at least \(1-\vartheta\).  They may depend on \(C_{\mathrm{lev}},\xi,\vartheta\).
\end{proposition}

The modified profile yields entropy dissipation and the higher logarithmic
moments after the reduction by level sets.  The classical profile yields a
support profile and hence the heat kernel bound.  We prove the profiles
in three ranges: \(L\le n^\delta\) by the argument with two observations;
\(n^\delta<L\le\lambda n\) by the proportional branch of the source
path; and \(\lambda n\le L\le C_{\mathrm{lev}}n\) by the product endpoint
and the estimate for a reweighted mean below, for a small fixed \(\lambda\)
chosen later.

\begin{lemma}[Entropy between two observation scales]\label{warm:small-profile}
Fix \(C_{\mathrm{lev}}<\infty\).  For every fixed \(\kappa>0\), there are tight \(K_0\) and
\(\delta>0\) such that every \(f\) satisfying
\eqref{warm:restricted-class} with \(L\le n^\delta\) satisfies
\begin{equation}
 \Tn\,\frac{\cE(f)}{N}\ge c_\kappa L^{5/3-\kappa}.
                                                               \label{warm:small-profile-bound}
\end{equation}
\end{lemma}

\begin{proof}
Fix first one instance of Proposition~\ref{prop:edge-estimates}, with
\(\epsilon_{\rm sp}\le1\), \(K\ge1\), and a variance-remainder exponent
\(0<a\le1\), decreasing the exponent if necessary.  Its constants are fixed on a
disorder event of arbitrarily high limiting probability.  Choose
\(M>1\) sufficiently large that
\[
 Ma>1,\qquad \frac{2M}{3}>\frac53-\kappa .
\]
We retain this variance estimate when imposing further cap tolerances
below; intersection with finitely many additional events does not change
its already chosen exponent \(a\).  Choose \(\delta>0\) so small that
\(M\delta/3<10^{-4}\).  This leaves every fixed factor buffer about the
scales below within \(r_n=n^{-1/3+10^{-4}}\) for large \(n\).
Use the same fixed small \(\gamma>0\) as in the covariance iteration
of Section~\ref{sec:edge-estimates}.  Choose the bridge tolerance
\(\epsilon>0\) after \(M\), with
\(2(M-1)\epsilon/3<\kappa\).

For a small fixed \(c_0>0\), to be chosen using only the projection-tail
constants, define
\[
 k_{p_1}=np_1^3=c_0L,\qquad k_R=nR^3=L^M.
\]
Increasing \(K_0\) puts these scales above the threshold of
Proposition~\ref{prop:edge-estimates}.  Suppose, for a contradiction, that
\(\Tn\cE(f)/N<L^{5/3-\kappa}\).  If
\(F_R(h)=\mu_{R,h}[f]\) and
\(V_R=Q_\mu[\Var(f\mid h_R)]\), that proposition gives
\begin{equation}
 \frac{V_R}{N}
 \le C L^{5/3-\kappa-2M/3}
       +C\exp\{L-cL^{Ma}\}=o_{K_0\to\infty}(1).
                                                               \label{warm:small-variance}
\end{equation}
The estimate is uniform in \(n,L,f\) in the present range, because
\(R^2\Tn=L^{2M/3}\).

Consider the joint ordinary law \(Q_\mu\) of the spin and its observation
path through \(R\).  Change this law to \(Q_{\rm spin}\) or \(Q_{\rm obs}\)
by the respective densities
\[
 \Lambda^{\rm spin}=\frac{f^2}{N},
 \qquad
 \Lambda^{\rm obs}=\frac{F_R(h_R)^2}{A},
 \qquad
 A=Q_\mu[F_R^2]=N-V_R .
\]
For \(K_0\) large, \(A\ge N/2\), so both densities are at most \(2e^L\).
For any of these laws, \(Q|_p\) denotes restriction to
\(\sigma(h_q:0\le q\le p)\), always excluding
the spin.  Define, for \(p<R\),
\begin{equation}\label{warm:residual-observation-entropy}
 H(p)=D(Q_{\rm obs}\Vert Q_\mu)
       -D(Q_{\rm obs}|_p\Vert Q_\mu|_p).
\end{equation}
The first entropy is on the full joint law.  Since its likelihood is
measurable from \(h_R\), it equals the relative entropy on the terminal
observation sigma-field.  Both terms are finite by the likelihood bound,
and the conditional entropy chain rule identifies \(H(p)\) with the
expected entropy of the terminal observation tilt remaining after the
prefix through \(p\).
To compare the two tilts, note that \(Q_\mu[fF_R]=A\), and hence
\[
 Q_\mu\left[\left(\frac f{\sqrt N}
                  -\frac{F_R}{\sqrt A}\right)^2\right]
 =2\left(1-\sqrt{\frac AN}\right)\le\frac{2V_R}{N}.
\]
Cauchy--Schwarz bounds the \(L^1(Q_\mu)\) distance between the two
densities by \(2\sqrt{2V_R/N}\).  This bound persists after projection
to any part of the observation path.  If two projected densities
\(u,v\le2e^L\) have \(L^1\) distance \(d\), then
\begin{equation}
 \left|Q_\mu[u\log u]-Q_\mu[v\log v]\right|
 \le (1+L+\log2)d+C.                                      \label{warm:entropy-continuity}
\end{equation}
Indeed \(x\log x\) has derivative bounded by \(1+L+\log2\) on
\([1,2e^L]\), while its absolute value is bounded on \([0,1]\).
It follows from \eqref{warm:small-variance} that the two relative
entropies, both before and after any projection, differ by \(o(L)\),
uniformly as \(K_0\) increases.

We bound the information revealed at \(p_1\) under the spin tilt.  The
field \(h_{p_1}\) is a sufficient statistic for the path through that
scale.  Let \(Q_0\) be the output of the same Gaussian channel with
zero signal.  Convexity of relative entropy and evenness of \(\mu\)
give
\[
 D(Q_{\rm spin}|_{p_1}\Vert Q_0)
       \le\frac12\frac{\mu[f^2X^{\mathsf T}B_{p_1}X]}N,
 \qquad
 \log\frac{\dd(Q_\mu|_{p_1})}{\dd Q_0}(h)
       \ge-\frac12\mu[X^{\mathsf T}B_{p_1}X].
\]
These formulas are on the range of \(B_{p_1}\), so zero precision
directions cause no difficulty.  Subtraction bounds
\(D(Q_{\rm spin}|_{p_1}\Vert Q_\mu|_{p_1})\) by half the sum of the two
quadratic moments.  The spectral bound \(d_1\ge-\epsilon_{\rm sp}p_1^2\)
gives \(B_{p_1}\preceq(1+\epsilon_{\rm sp})p_1^2P_{p_1}
\preceq2p_1^2P_{p_1}\).  The projection
tail in Proposition~\ref{prop:edge-estimates}, integrated against a density at most
\(e^L\), gives
\[
 \frac{\mu[f^2\|P_{p_1}X\|^2]}{NM_{p_1}^2}
       +\frac{\mu[\|P_{p_1}X\|^2]}{M_{p_1}^2}
 \le C\left(1+(L/k_{p_1})^{1/3}\right).
\]
For example, integrate
\(\min\{1,e^LCe^{-c k_{p_1}v^6}\}\) in the variable \(v^2\).
The information is therefore at most
\[
 Ck_{p_1}\left(1+(L/k_{p_1})^{1/3}\right)
   \le C(c_0+c_0^{2/3})L .
\]
This is the proportional-cap argument in
\cite[Equation~(19.12)]{CM}.  Choose \(c_0\) so that this last bound is
less than \(L/8\).  The full spin-tilt entropy is at least \(L/2\).
The two continuity comparisons in \eqref{warm:entropy-continuity},
combined into one uniform \(o(L)\), therefore give
\begin{equation}\label{warm:residual-lower}
 \begin{split}
 H(p_1)&\ge D(Q_{\rm spin}\Vert Q_\mu)
       -D(Q_{\rm spin}|_{p_1}\Vert Q_\mu|_{p_1})-o(L)\\
       &\ge L/2-L/8-o(L)\ge cL.
 \end{split}
\end{equation}

We follow that remaining entropy up to a fixed fraction of \(R\).
For \(p<R\), let \(\mathsf T_{p,R}^h\) be the ordinary conditional
law of \(h_R\) given \(h_p=h\).  It has density proportional to
\[
 e^{\phi_R(y)}\,\mathcal N(h,B_R-B_p)(\dd y),
\]
where the Gaussian is restricted to the range of its covariance.
Let \(\widetilde{\mathsf T}_{p,R}^h\) be its tilt by \(F_R(y)^2\), and
put
\[
 A_p(h)=\mathsf T_{p,R}^h[F_R^2],\qquad
 l_p(h)=D(\widetilde{\mathsf T}_{p,R}^h\Vert\mathsf T_{p,R}^h).
\]
Observation sufficiency gives the prefix density
\(\dd(Q_{\rm obs}|_p)/\dd(Q_\mu|_p)=A_p(h_p)/A\).  The conditional
chain rule in \eqref{warm:residual-observation-entropy} now reads
\begin{equation}\label{warm:residual-transition-entropy}
 H(p)=\frac1A Q_\mu\!\left[
       \Ent_{\mathsf T_{p,R}^{h_p}}(F_R^2)\right]
       =Q_{\rm obs}[l_p(h_p)].
\end{equation}

Write \(\Delta=B_R-B_p\).  On the currently active subspace
\(\cH_p\), \(\Delta=(R^2-p^2)I\).  Differentiating its Gaussian
convolution therefore gives
\[
 \nabla_{\cH_p}\log A_p(h)=\delta m_h,\qquad
 \delta m_h=\widetilde{\mathsf T}_{p,R}^h[S]
                    -\mathsf T_{p,R}^h[S],\qquad
 S=\frac{P_p(h_R-h)}{R^2-p^2}.
\]
In the parameter \(p^2\), the likelihood martingale \(A_p(h_p)/A\)
has innovation coefficient \((A_p/A)\delta m_h\).  The Gaussian
likelihood formula for its relative entropy and the chain rule yield
\begin{equation}
 -\frac{\dd H}{\dd\log p}
      =p^2Q_{\rm obs}[|\delta m_h|^2] .                   \label{warm:bridge-loss}
\end{equation}
This identity holds almost everywhere and in integrated form across
activation points.  One direct justification applies It\^o's formula
to \((u+\epsilon)\log(u+\epsilon)\) of the bounded likelihood,
then lets \(\epsilon\downarrow0\); zero-rate directions are omitted.
This is the observation entropy identity
\cite[Lemma~19.2, Equation~(19.6)]{CM}, applied here to a terminal
observation tilt instead of a spin tilt.

The ordinary log moment generating function of \(S\), for
\(v\in\cH_p\), is exactly
\[
 \log\mathsf T_{p,R}^h[e^{\ip vS}]
   =\phi_p(h+v)-\phi_p(h)+\frac{|v|^2}{2(R^2-p^2)} .
\]
Suppose that the Hessian of \(\phi_p\) in this subspace is at most
\(C_{\rm curv}p^{-2}I\) on all shifts of length at most
\(p\sqrt{2l_p(h)/C_{\rm curv}}\) from \(h\).  The entropy variational
inequality, tested against \(\ip vS\), then gives
\begin{equation}
 p^2|\delta m_h|^2
 \le2\left(C_{\rm curv}+\frac{p^2}{R^2-p^2}\right)l_p(h).
                                                               \label{warm:bridge-mean}
\end{equation}
To check the finite shift radius, the centered log moment generating
function is at most
\(\frac12(C_{\rm curv}p^{-2}+(R^2-p^2)^{-1})|v|^2\).
Testing at length \(\sqrt{2l_p/(C_{\rm curv}p^{-2}+(R^2-p^2)^{-1})}\)
in the mean-shift direction proves the inequality by contradiction;
this length is no larger than the assumed radius.

We next account for observations where that curvature bound is not
available.  The entropy bound \(l_p\le-\log A_p\), together with
\(A\ge e^{-L}/2\), implies
\begin{equation}
 Q_{\rm obs}(l_p>L+\log2+z)\le e^{-z},\qquad z\ge0.
                                                               \label{warm:bridge-entropy-tail}
\end{equation}
Under the ordinary joint law,
\[
 S=P_pX+(R^2-p^2)^{-1/2}G_p,
\]
with an independent standard Gaussian \(G_p\) on \(\cH_p\).
For \(p\le\gamma R\), the projection tails of
Proposition~\ref{prop:edge-estimates} and Gaussian concentration give,
for each prescribed \(a_0>0\), a fixed \(V\) such that
\begin{equation}\label{warm:bridge-input-tails}
 \begin{split}
 Q_\mu\!\left[\frac{|S|^2}{M_p^2}
       \ind_{\{|S|>VM_p\}}\right]&\le Ce^{-a_0k_p},\\
 Q_\mu\{\|h_p\|>Vp^2M_p\}&\le Ce^{-a_0k_p}.
 \end{split}
\end{equation}
To see the common scale, write the observation noise using a standard
Gaussian \(G'_p\) on \(\cH_p\).  Since \(B_p\preceq2p^2P_p\),
\[
 \frac{\|h_p\|}{p^2M_p}
 \le2\frac{\|P_pX\|}{M_p}+\sqrt2\frac{\|G'_p\|}{\sqrt{k_p}},
 \qquad
 \frac{|S|}{M_p}
 \le\frac{\|P_pX\|}{M_p}
       +\frac p{\sqrt{R^2-p^2}}\frac{\|G_p\|}{\sqrt{k_p}}.
\]
The last coefficient is bounded on this range and
\(\dim\cH_p\le Ck_p\).  Integration of the same tails allows any
fixed polynomial moment as well.  For a prefix event \(E\), conditional Jensen and
the full and prefix likelihood bounds at most \(2e^L\) give
\begin{equation}
\begin{split}
 p^2Q_{\rm obs}[|\delta m_h|^2\ind_E]
 \le Ck_p\bigg(
 V^2Q_{\rm obs}(E)
 +e^L Q_\mu\!\left[
       \frac{|S|^2}{M_p^2}\ind_{\{|S|>VM_p\}}\right]\bigg).
\end{split}                                                     \label{warm:discarded-bridge}
\end{equation}
Indeed Jensen bounds \(|\delta m_h|^2\) by twice the sum of the
second moments under the two conditional laws.  For their parts
\(|S|\le VM_p\) use the prefix probability of \(E\); for the
remaining parts use respectively the full likelihood \(F_R^2/A\)
and the prefix likelihood \(A_p/A\).  Finally \(p^2M_p^2=k_p\).

Choose \(p_2\) by \(k_{p_2}=C_2k_{p_1}\), where \(C_2\) will be
fixed below.  Once it is fixed, \(p_2<\gamma R\) for \(K_0\)
sufficiently large.  Fix \(A_0>0\), choose a tail rate
\(a_0>c_0^{-1}\) with a fixed margin, and then choose \(V\) in
\eqref{warm:bridge-input-tails}.  On \([p_1,p_2]\) retain
\(\|h_p\|\le Vp^2M_p\) and \(l_p\le L+\log2+A_0k_p\).
The trial shifts in \eqref{warm:bridge-mean}, in units \(p^2M_p\),
then have length at most
\(\sqrt{2(c_0^{-1}+A_0+1)}\), after enlarging the ball-curvature
constant to be at least one.  Choose a fixed ball containing the retained
fields and these shifts before requesting its curvature constant from
Proposition~\ref{prop:edge-estimates}.  Equations
\eqref{warm:bridge-entropy-tail}--\eqref{warm:discarded-bridge}
show that discarded observations contribute at most \(Ck_pe^{-c k_p}\).
Thus \eqref{warm:bridge-loss}--\eqref{warm:bridge-mean} imply
\[
 \frac{\dd H}{\dd\log p}\ge-C_1H-Ck_pe^{-c k_p}
       \qquad(p_1\le p\le p_2).
\]
The logarithmic length of this interval is fixed.

For \([p_2,\gamma R]\), make a fresh auxiliary cap and extension
request with the same \(\gamma\).  Choose \(\eta\ll\gamma^2\), then
\(\epsilon_{\rm aux}\le\epsilon\) sufficiently small compared with
\(\eta\).  Reapply the construction of Lemma~\ref{edge:annular-extension}
with these choices, and write \(s_0\) for its auxiliary width.
Let \(\mathcal G_{\rm ext}=\mathcal G_n^{\rm aux,ext}
\cap\mathcal G_n^{\rm FUA}\) be the resulting event for the required
static tolerances, radii, threshold, and scale buffers.
Lemma~\ref{edge:cube-caps} gives \(\mathcal G_n^{\rm aux,ext}\) with
arbitrarily high limiting ordinary probability, and
\(\mathcal G_n^{\rm FUA}\) costs \(o(1)\).  Assign a further fraction
of the failure probability and intersect with the retained variance
event; its exponent \(a\) and constants are unchanged.

On the full annulus \(|\alpha_p(h)|\le s_0p^2\), this request gives
\(K_p(h)\preceq(1+\epsilon_{\rm aux})p^{-2}I
\preceq(1+\epsilon)p^{-2}I\).  The same annulus specifies admissible
child fields for the extension.  Choose \(0<\zeta<s_0\) and a fixed
\(b_{\rm rad}>0\) so that shifts of length at most
\(b_{\rm rad}p^2M_p\) from \(|\alpha_p(h)|\le\zeta p^2\) stay in the
full annulus.  After \(\zeta\) is fixed, impose its ordinary annulus
estimate from Lemma~\ref{edge:path-annuli} as one further fixed request
and intersection.  Thus the ordinary failure probability is
\(Ce^{-a_\zeta k_p}\), with a fixed \(a_\zeta>0\); only this probability
bound is used from the further request.  The trial shifts fit whenever
\(l_p\le b k_p\), where \(b=(1+\epsilon)b_{\rm rad}^2/2\).
Choose \(C_2\) so large that
\[
 \frac{L}{k_p}\le\frac1{C_2c_0}
       <\frac14\min\{b,a_\zeta\}\qquad(p\ge p_2).
\]
The likelihood bound and \eqref{warm:bridge-entropy-tail} give
exponentially small tilted probabilities for both discarded events.
Using \eqref{warm:discarded-bridge} with a further fixed \(S\)-tail
threshold proves
\[
 \frac{\dd H}{\dd\log p}
 \ge-\left(2(1+\epsilon)+\frac{2p^2}{R^2-p^2}\right)H
        -Ck_pe^{-c k_p}.
\]
The extra coefficient has bounded integral,
\[
 \int_{p_2}^{\gamma R}\frac{2p^2}{R^2-p^2}\,\dd\log p
       \le-\log(1-\gamma^2).
\]
In both intervals the integrating factor charges a fixed power of
\(p/p_1\), whereas \(k_p=np^3\) grows as its cube.  Consequently the
weighted errors are bounded by a polynomial in \(k_{p_1}\) times
\((e^{-c k_{p_1}}+e^{-c k_{p_2}})\).  They are \(o(L)\) after
increasing \(K_0\).  The first interval costs a fixed multiplicative
factor.  We obtain
\begin{equation}
 H(\gamma R)\ge cL(p_1/R)^{2+2\epsilon}.                  \label{warm:retained-small}
\end{equation}

It remains to bound this entropy by the energy of \(f\).  For a child
field satisfying \(|\alpha_{\gamma R}(h)|\le s_0(\gamma R)^2\), the
fresh request to Lemma~\ref{edge:annular-extension} supplies an extension
of the parent potential with upper Hessian \((1+\eta)R^{-2}I\).  Put
\[
 \mathcal A_R=\{|\alpha_R|\le s_0R^2/2\},\qquad
 \mathcal A_R^+=\{|\alpha_R|<3s_0R^2/4\}.
\]
The extension agrees with the true potential on \(\mathcal A_R^+\).
The distance from \(\mathcal A_R\) to \((\mathcal A_R^+)^c\) is at least
\(c_{\rm buf}R^2M_R\).  Both transition laws assign
\(\mathcal A_R^c\) probability \(Ce^{-c k_R}\), also with fixed
polynomial moments of \(\|h_R\|/(R^2M_R)\).
The precision increment satisfies
\(\Delta\preceq(1-\gamma^2)R^2I\).  On its range the extended
transition, denoted by \(\overline{\mathsf T}\), has negative log-density Hessian at least
\[
 \Delta^{-1}-(1+\eta)R^{-2}I\succeq c_\gamma R^{-2}I .
\]
The Bakry--\'Emery logarithmic Sobolev criterion
\cite[Corollary~5.7.2]{BakryGentilLedoux2014} gives
\[
\Ent_{\overline{\mathsf T}}(g^2)
\le C_\gamma R^2\overline{\mathsf T}[|\nabla g|^2].
\]
It applies to the semiconcave extension by mollification.

Choose a smooth cutoff \(\chi\) equal to one on \(\mathcal A_R\)
and supported in \(\mathcal A_R^+\).  The fixed buffer and the
radial saddle derivative in Proposition~\ref{prop:edge-estimates} give
\(|\nabla\chi|\le C/(R^2M_R)\).  Apply the preceding inequality to
\(\chi F_R\).  Its cutoff-gradient term is at most \(Ce^{-c k_R}\),
because \(R^2|\nabla\chi|^2\le C/k_R\) and the transition tail has
that exponential rate.  On agreement the ratio of the two transition
densities is constant; their normalization and their exponentially small
tails make that constant \(1+O(e^{-c k_R})\).  To compare the entropies,
use
\(\Ent_{\mathsf T}(g)=\mathsf T[g\log g]-\mathsf T[g]\log\mathsf T[g]\):
\(x\log x\) is bounded and H\"older continuous on \([0,1]\), and
\(0\le F_R,\chi\le1\).  The comparison, with a possibly smaller
positive exponent, proves
\begin{equation}
 \Ent_{\mathsf T_{\,\gamma R,R}^{h}}(F_R^2)
 \le C_\gamma R^2\mathsf T_{\,\gamma R,R}^{h}
                   [|\nabla_{\cH_R}F_R|^2]+Ce^{-c k_R}.
                                                               \label{warm:transition-lsi}
\end{equation}
Failure of the child condition is contained in the ordinary
\(\zeta\)-annulus failure, since \(\zeta<s_0\).  Its probability is
\(Ce^{-c k_R}\) after changing \(c\), since
\(k_{\gamma R}=\gamma^3k_R\).
On it the conditional entropy of \(F_R^2\le1\) is at most \(1/e\).
Average \eqref{warm:transition-lsi} under the ordinary child law and
use \eqref{warm:residual-transition-entropy}.  The tower property,
\(A\ge N/2\), and the two inequalities of
\eqref{edge:variance-estimates} at \(R\) give
\begin{equation}
 \begin{split}
 H(\gamma R)
 &\le\frac{C_\gamma R^2}{A}
       Q_\mu\!\left[\|\nabla_{\cH_R}F_R(h_R)\|^2\right]
       +\frac{Ce^{-c k_R}}A\\
 &\le C\frac{V_R}{N}+C\exp\{L-c k_R^a\}\\
 &\le C R^{-2}\frac{\cE(f)}N+C\exp\{L-c k_R^a\}.
 \end{split}                                                \label{warm:small-upper}
\end{equation}
Here \(0<a\le1\) allows the transition error to be included in the
displayed remainder, and \(k_R=L^M\) makes it \(C e^{L-cL^{Ma}}\).
The error is negligible compared with
\eqref{warm:retained-small}, even if the latter is a negative power
of \(L\).  Comparing the two bounds and using
\(\Tn p_1^2=(c_0L)^{2/3}\) gives
\[
 \Tn\,\frac{\cE(f)}N
 \ge c\,\Tn R^2L(p_1/R)^{2+2\epsilon}
 =c' L^{5/3-2(M-1)\epsilon/3}.
\]
The choice \(2(M-1)\epsilon/3<\kappa\) contradicts the assumed upper
bound for \(K_0\) large.  All ball radii, tail rates, annulus
tolerances and factor buffers have been fixed before this last
increase of \(K_0\) and before \(n\) tends to infinity.
\end{proof}

\subsection{The product endpoint and larger entropy levels}
\label{warm:large-section}

We record first two entropy bounds at a product posterior.  Let
\(\nu=f^2\mu/N\).  For each path used below, let \(Q_\mu,Q_\nu\)
denote the corresponding joint laws of the initial spin and its
observations, started from these two spin laws.  Expectations of
observation-only functions use their observation marginals, and
\(Q|_u\) again denotes the observation prefix, excluding the spin.
At observation time \(u\), write \(\mu_u,\nu_u\) for their spin
posteriors and \(N_u=\mu_u[f^2]\).  Use the posterior heat-bath form and
the weighted Euclidean norm
\[
 \cE_{\mu_u}(f)=\sum_i\mu_u[\Var_{\mu_u}(f\mid X_{-i})],
 \qquad \|v\|_A^2=v^{\mathsf T}Av.
\]
Bayes' formula and the observation
identities of \cite[Lemma~19.2, Equations~(19.6)--(19.8)]{CM} are
\begin{align}
 \frac{\dd(Q_\nu|_u)}{\dd(Q_\mu|_u)}
       &=\frac{N_u}{N},\qquad \nu_u=\frac{f^2\mu_u}{N_u},
                                                               \label{warm:posterior-likelihood}\\
 -\dot H(u)
       &=\frac12 Q_\nu\!\left[
             \|\nu_u[X]-\mu_u[X]\|_{\dot B_u}^{\,2}\right],
 \qquad H(u)=Q_\nu[D(\nu_u\Vert\mu_u)],
                                                               \label{warm:spin-loss}\\
 Q_\nu\!\left[\frac{\cE_{\mu_u}(f)}{N_u}\right]
       &=\frac1N Q_\mu[\cE_{\mu_u}(f)]
       \le\frac{\cE(f)}N .                                \label{warm:conditioned-energy}
\end{align}
The posterior \(\nu_u\) is defined where \(N_u>0\); the remaining
observation paths have \(Q_\nu\)-mass zero and contribute zero to
these averaged formulas.
The loss identity has the corresponding stopped-path form, with the
information in precision jumps added.  The last inequality is
conditional-variance contraction, site by site, under the ordinary
joint spin-observation law.  The form at every posterior uses the
same rate-one-per-site normalization as \(\cE\).

\begin{lemma}[Product-posterior entropy bounds]\label{warm:product}
Fix the at most two product endpoints used below.  Each has precision
\(t'I+W\) for a fixed \(t'>3\), chosen from the fixed path parameters
before \(f\).  On the common disorder event of
\cite[Equation~(19.10)]{CM} for these endpoints, including its
polynomial-moment field-tail control, the tests
\eqref{warm:restricted-class} satisfy
\begin{equation}
 H_{\rm end}\le
 C\sqrt{\log n+L}\,\frac{\cE(f)}N+O(n^{-3}),
 \qquad
 H_{\rm end}\le\frac{\cI(f^2)}N .                         \label{warm:product-bounds}
\end{equation}
The constants in the first inequality, including its remainder, may
depend on these endpoints, the fixed path parameters, and
\(C_{\mathrm{lev}}\), but not on \(n\) or \(f\).
The second inequality holds for every observation channel whose
posterior at the endpoint is a product measure; it does not require
a field-tail estimate.
\end{lemma}

\begin{proof}
The first inequality is the product calculation of
\cite[Equation~(19.18)]{CM}, applied at each chosen endpoint.
Its precision is \(t'I+W\) for the corresponding fixed \(t'>3\).
The posterior interaction is then a scalar diagonal, hence its spin
law is a product with fields \(h_i\).  Tensorizing the two-point
log-Sobolev bound gives
\[
 D(\nu_h\Vert\mu_h)
 \le C(1+\|h\|_\infty)\frac{\cE_{\mu_h}(f)}{N_h}.
\]
The field tail of \cite[Equation~(19.10)]{CM}, on the common event for
the chosen endpoints, restricts
\(\|h\|_\infty\le C\sqrt{\log n+L}\).  Its ordinary complement has
probability at most \(e^{-2L}n^{-10}\), with polynomial moments;
after the likelihood bound \(N_h/N\le e^L\) its entropy contribution
is \(O(n^{-3})\).  Equation \eqref{warm:conditioned-energy} gives the
displayed bound.

For the second inequality, entropy tensorization at a product posterior
and Jensen's inequality for the logarithm give, for positive \(z\),
\[
 \Ent_{\mu_h}(z)
 \le\sum_i\mu_h[\Ent_{\mu_h(\cdot\mid X_{-i})}(z)]
 \le\sum_i\mu_h[\Cov_{\mu_h(\cdot\mid X_{-i})}(z,\log z)]
 =\cI_{\mu_h}(z).
\]
In fact \(\Ent_\pi(z)\le\Cov_\pi(z,\log z)\) follows from
\(\pi[\log z]\le\log\pi[z]\).  The posterior entropy formula is
\[
 H_{\rm end}=\frac1N Q_\mu[\Ent_{\mu_h}(f^2)].
\]
It remains to compare the averaged modified form to its initial value.
For a fixed site and fixed \(x_{-i}\), let
\[
 a_h=\mu_h(X_{-i}=x_{-i},X_i=+1),\qquad
 b_h=\mu_h(X_{-i}=x_{-i},X_i=-1)
\]
be the two joint posterior masses.  Their contribution to
\(\cI_{\mu_h}(z)\) is
\[
 \frac{a_hb_h}{a_h+b_h}
       (z_+-z_-)(\log z_+-\log z_-).
\]
The second factor is nonnegative and independent of the observation.
The function \(c(a,b)=ab/(a+b)\), extended by zero on the coordinate
axes, is concave: in the positive quadrant its Hessian is
\(-2(b,-a)(b,-a)^{\mathsf T}/(a+b)^3\).
Since \(Q_\mu[a_h]\) and \(Q_\mu[b_h]\) are the corresponding
initial masses, Jensen's inequality proves
\(Q_\mu[\cI_{\mu_h}(z)]\le\cI_\mu(z)\).
Take \(z=f^2+\epsilon\) and let \(\epsilon\downarrow0\) to include
zeros.  This proves the second bound.
\end{proof}

At entropy of order \(n\), the observation path must be stopped before
its field leaves a region where the posterior mean is controlled.
The following strengthening of the square-density estimate in
\cite[Proposition~7.3]{SG} is what permits the necessary energy range.
We display the stability condition used at subcritical posterior
interactions; the finite word estimates and vector recipes of \cite{SG}
are used in their stated form.

Let \(J_*=\beta W_{\rm off}\), \(0<\beta<1\), and \(j=\beta^2\).
For \(y,h\in\R^n\) put
\[
\begin{gathered}
 m(y)=\tanh y,\qquad q(y)=\frac{\|m(y)\|^2}{n},\qquad
 V(y)=\diag(1-m_i(y)^2),\\
 F_h(y)=y+\bigl(j(1-q(y))I-J_*\bigr)m(y)-h .
\end{gathered}
\]
The stability condition of \cite[Equation~(17.2)]{CM}, with fixed
positive \(b,\eta,\epsilon_A,\rho_0\), is the implication
\begin{equation}
\begin{gathered}
 n^{-1/2}\|F_h(y)\|\le2\rho_0,\qquad 0\preceq A\preceq(1+\eta)I,\\
 A\ \hbox{diagonal},\qquad
 n^{-1/2}\|A-V(y)\|_{\rm F}\le\epsilon_A\\
\Longrightarrow\quad
 I+A^{1/2}\left(j(1-q(y))I-J_*
       -\frac{2j}{n}m(y)m(y)^{\mathsf T}\right)A^{1/2}\succeq bI .
\end{gathered}
                                                               \label{warm:stability}
\end{equation}
for every \(y,A\).  The ordinary word diagnostics below are the
simultaneous estimates of \cite[Lemma~3.1]{SG}: for a prescribed
finite length they bound the operator norm, the off-diagonal entrywise
\(\ell^4\) norm, and the Euclidean norm of the diagonal error from
the noncrossing prediction for products of \(J_*\) and bounded
diagonal matrices.  The ordinary vector recipes are those of
\cite[Definition~6.2]{SG}; their square-density argument uses no
inverse word.

\begin{lemma}[Square-root energy in a reweighted posterior mean]
\label{warm:posterior-mean}
Fix an operator-norm bound for \(J_*\) and the parameters in
\eqref{warm:stability}, with \(j(1+\eta)^2<1\).
For every sufficiently small fixed \(\epsilon>0\), impose the
ordinary word diagnostics of a finite length depending on \(\epsilon\).
For any field \(h\) satisfying \eqref{warm:stability}, taking
\(A=V(y)\) verifies the hypothesis of \cite[Lemma~5.7]{SG} with
residual threshold \(2\rho_0\sqrt n\) and lower bound \(b\).
Let \(y_*\) be its unique root of \(F_h\), and set
\(m_*=\tanh y_*\).  This root existence is the source conclusion;
the square-root energy dependence below is the new estimate.  Define
\[
 \mu_h(\sigma)\ \propto\
   \exp\{\tfrac12\sigma^{\mathsf T}J_*\sigma+h^{\mathsf T}\sigma\},
 \qquad
 \mathcal D_h(f):=\cE_{\mu_h}(f)
       =\sum_i\mu_h[\Var(f\mid\sigma_{-i})].
\]
For every real \(f\) with \(N_h=\mu_h[f^2]>0\) and
\(\nu_h=f^2\mu_h/N_h\),
\begin{equation}
 \frac{\|\nu_h[\sigma]-m_*\|}{\sqrt n}
 \le C_{\rm lead}\epsilon+
       \frac{C_\epsilon}{\sqrt n}
          \left(1+\sqrt{\frac{\mathcal D_h(f)}{N_h}}\right).
                                                               \label{warm:sqrt-posterior}
\end{equation}
The constants are uniform in \(n,h,f\) under these hypotheses.
\(C_{\rm lead}\) depends only on the fixed stability and norm bounds;
the finite diagnostic length and \(C_\epsilon\) may depend on
\(\epsilon\).  The statement is simultaneous on a fixed compact
subinterval of \((0,1)\) carrying a common strict margin, as in the
last paragraph of the proof of \cite[Proposition~17.1]{CM}.
\end{lemma}

\begin{proof}
We prove the stronger square-density residual estimate that replaces
the linear energy dependence in \cite[Lemma~6.5]{SG}.  Use the
zero-diagonal interaction \(J_*\), so the cube coordinates in this
proof are spins, not the scaled edge coordinates.  For a spin flip
\(\sigma^{(i)}\), put
\[
 d_i g(\sigma)=\frac{g(\sigma)-g(\sigma^{(i)})}{2\sigma_i}.
\]
The primary iterates of \cite[Section~6]{SG} are
\[
\begin{gathered}
 m^0=\sigma,\qquad h^1=h+J_*\sigma,\qquad
 m^l=\tanh h^l,\qquad b_l=1-\|m^l\|^2/n,\\
 h^{l+1}=h+J_*m^l-jb_lm^{l-1},\qquad
 R_l=m^{l-1}-m^l .
\end{gathered}
\]
For a fixed ordinary recipe admissible from level \(l\), let \(U\)
be a terminal vector in that recipe.  We claim that, whenever
\(\sup_\sigma(|H_0|+\|dH_0\|)\le C_0\),
\begin{equation}
 \left|\mu_h[f^2H_0R_l^{\mathsf T}U]\right|
 \le C\left(N_h+\sqrt{N_h\mathcal D_h(f)}\right).
                                                               \label{warm:sqrt-residual}
\end{equation}
The constant has the dependence on the fixed recipe, depth, coefficient
bounds and finite word diagnostics allowed in the source lemma, and
is independent of \(f\).  The residual claim itself does not require
root stability.

Here are the base case and the induction, including the point where
the improvement occurs.  The pointwise estimate
\cite[Lemma~6.4, Equation~(6.20)]{SG} gives
\[
 \|U\|+\sum_i|d_iU_i|\le C .
\]
It also gives the same bound for
\(\widetilde U=H_0U\): indeed
\[
 d_i(H_0U_i)=H_0\,d_iU_i+U_i(\sigma^{(i)})d_iH_0,
 \qquad
 \|(U_i(\sigma^{(i)}))_i\|
       \le\|U\|+2\|\diag\nabla U\|_2\le C,
\]
and Cauchy--Schwarz bounds the sum of the added diagonal derivatives.
Write \(C_U\) for a bound on
\(\|\widetilde U\|+\sum_i|d_i\widetilde U_i|\).

Let \(p_i^\pm=\mu_h(\sigma_i=\pm1\mid\sigma_{-i})\).
Then \(m_i^1=p_i^+-p_i^-\) and the conditional variance is
\(v_i^1=4p_i^+p_i^-\), in particular \(v_i^1\le4p_i^\pm\)
for each sign.  Since a function of one spin is affine with slope
\(d_if\), the form is
\[
 \mathcal D_h(f)=\mu_h\sum_i v_i^1(d_if)^2 .
\]
The exact one-site integration by parts
\cite[Equation~(6.31)]{SG} gives
\begin{align*}
 \mu_h[f^2R_1^{\mathsf T}\widetilde U]
 &=\mu_h\sum_i v_i^1d_i(f^2\widetilde U_i)\\
 &=\mu_h\sum_i v_i^1 f^2d_i\widetilde U_i
   +\mu_h\sum_i v_i^1(f_{i,+}+f_{i,-})d_if\,
                                  \widetilde U_i(\sigma^{(i)}).
\end{align*}
Here \(f_{i,\pm}=f(\sigma^{i,\pm})\), and
\(d_i(f^2)=(f_{i,+}+f_{i,-})d_if\), also for signed \(f\).
The first term is at most \(C_UN_h\) in absolute value.
For the second set
\[
 M_i(\sigma_{-i})=\max_\pm|\widetilde U_i(\sigma^{i,\pm})|.
\]
At each configuration,
\[
 \left(\sum_iM_i(\sigma_{-i})^2\right)^{1/2}
 \le\|\widetilde U(\sigma)\|
       +2\|\diag\nabla\widetilde U(\sigma)\|_2\le 2C_U.
\]
Weighted Cauchy--Schwarz bounds that second term by
\[
 \mathcal D_h(f)^{1/2}
 \left(\mu_h\sum_i v_i^1(f_{i,+}+f_{i,-})^2M_i^2\right)^{1/2}.
\]
The two separate bounds \(v_i^1\le4p_i^\pm\) imply
\[
 v_i^1(f_{i,+}+f_{i,-})^2
       \le8(p_i^+f_{i,+}^2+p_i^-f_{i,-}^2).
\]
Since \(M_i\) is independent of \(\sigma_i\), conditioning term by
term gives
\[
 \mu_h\sum_i v_i^1(f_{i,+}+f_{i,-})^2M_i^2
       \le8\mu_h\!\left[f^2\sum_iM_i^2\right]\le32C_U^2N_h .
\]
This proves \eqref{warm:sqrt-residual} for \(l=1\).

The finite induction in the source preserves this new right side.
For clarity, let \(U\) be admissible from level \(l+1\), put
\(a_l=j(b_l-b_{l-1})\), and define
\[
 U_i'=\Phi(h_i^l;h_i^{l+1},a_l)U_i,\qquad
 \Phi(z;v,a)=\int_0^1
 e^{a(1-u^2)/2}\frac{\cosh(u(z-v))}{\cosh z\cosh v}\,\dd u .
\]
Here \(|a_l|\le j\).  The integral formula bounds \(\Phi\) and every
needed fixed-order derivative uniformly in its two real field arguments:
the hyperbolic derivatives introduce bounded hyperbolic-tangent factors,
and the exponential factor and its derivatives are bounded for
\(|a|\le j\).  Together with the explicit level indices, these bounds
make \(U'\) admissible from level \(l\).  The scalar identity in
\cite[Section~6]{SG} then gives the residual expansion below.  Adjoin
the auxiliary field
\[
 y'=J_*U'-j\sum_q\theta_qm^{q-1}
              -j\sum_b\omega_bw^b,\qquad
 \theta_q=\frac1n\sum_i\partial_qU'_i,\quad
 \omega_b=\frac1n\sum_i\partial_{y^b}U'_i .
\]
Here \(w^b,y^b\) are the finitely many earlier auxiliary sources and
fields of the recipe; the partials differentiate their explicitly
displayed local arguments, holding parameters fixed, as in
\cite[Definition~6.2]{SG}.  Equation~(6.33) of that source is
\begin{align*}
 R_{l+1}^{\mathsf T}U
 ={}&R_l^{\mathsf T}y'+a_lR_l^{\mathsf T}U'
       -jb_{l-1}R_{l-1}^{\mathsf T}U'
       +j\sum_b\omega_bR_l^{\mathsf T}w^b\\
 &+j\sum_{q\ge l}\theta_qR_l^{\mathsf T}m^{q-1}
       -a_ln\theta_l ,
\end{align*}
with the \(R_{l-1}\) term absent when \(l=1\).  All first-line terms
are earlier residuals with ordinary admissible sources.  Their scalar
multipliers and their discrete-gradient norms are bounded by the
fixed-recipe bounds.  For \(q>l\) in the second line, use the source
\(m^{q-1}/\sqrt n\) and multiplier \(j\sqrt n\,\theta_q\);
\cite[Equation~(6.21)]{SG} gives the same boundedness for that
multiplier.  Finally the exact identity
\[
 R_l^{\mathsf T}m^{l-1}
       =n(b_l-b_{l-1})-R_l^{\mathsf T}m^l
\]
cancels the displayed \(-a_ln\theta_l\).  The remainder again uses
an earlier residual, the source \(m^l/\sqrt n\), and a bounded
multiplier.  A bounded multiplier \(\zeta\) with bounded \(\|d\zeta\|\)
can always be absorbed into \(H_0\), because
\[
 d_i(\zeta H_0)=\zeta(\sigma^{(i)})d_iH_0+H_0(\sigma)d_i\zeta .
\]
Every induction call therefore retains the original \(f,N_h,\mathcal D_h(f)\).
There are finitely many calls with a smaller residual index.  This
proves \eqref{warm:sqrt-residual} at every fixed depth.

We spell out its consequences for selection and the mean.  Put
\(S_l=R_l^{\mathsf T}m^l/\sqrt n\).
The primary derivative bounds in \cite[Lemma~7.1]{SG} give
\(|S_l|\le2\sqrt n\) and \(\|dS_l\|\le C_l\).  Apply
\eqref{warm:sqrt-residual} with \(U=m^l/\sqrt n\) and
\(H_0=S_l/\sqrt n\).  Division by \(N_h\sqrt n\) yields the
strengthened form of \cite[Equation~(7.1)]{SG}:
\begin{equation}
 \nu_h\!\left[\left(\frac{S_l}{\sqrt n}\right)^2\right]
 \le\frac{C_l}{\sqrt n}
       \left(1+\sqrt{\frac{\mathcal D_h(f)}{N_h}}\right).
                                                               \label{warm:sqrt-selection-moment}
\end{equation}
For a small fixed \(\rho>0\), choose a fixed depth \(d\) with
\(\lfloor d/2\rfloor\rho^2/4>2\), and then \(e_\rho>0\) with
\(2de_\rho<1\).  The telescoping identity
\[
 \frac{\|R_l\|^2}{n}=b_l-b_{l-1}-\frac{2S_l}{\sqrt n}
\]
shows that \(\max_{l\le d}|S_l|/\sqrt n\le e_\rho\) forces a
consecutive pair \(R_{k-1},R_k\) to have both norms at most
\(\rho\sqrt n/2\).  Otherwise the disjoint pairs would make the sum
of squared norms exceed \(2\), whereas telescoping bounds it by
\(1+2de_\rho<2\).

As in that source lemma, choose smooth \(C_k\in[0,1]\), equal to one
on this smaller pair region and supported where both norms are at
most \(\rho\sqrt n\).  The primary derivative bounds give
\(\|dC_k\|\le C_\rho/\sqrt n\).  Put
\[
 \widehat C_k=C_k\prod_{i=2}^{k-1}(1-C_i),\qquad
 D_{\rm sel}=1-\sum_{k=2}^d\widehat C_k=\prod_{k=2}^d(1-C_k).
\]
Positive deficit implies \(\max_{l\le d}|S_l|/\sqrt n>e_\rho\).
Chebyshev's inequality and \eqref{warm:sqrt-selection-moment} give
\begin{equation}
 \nu_h(D_{\rm sel}>0)
 \le\frac{C_\rho}{\sqrt n}
       \left(1+\sqrt{\frac{\mathcal D_h(f)}{N_h}}\right).
                                                               \label{warm:sqrt-deficit}
\end{equation}
The cutoffs are admissible bounded multipliers in
\eqref{warm:sqrt-residual}.

Choose \(\rho\le\epsilon\) with
\((\|J_*\|+4j)\rho<2\rho_0\).
On the support of \(C_k\), the recursion gives
\[
 F_h(h^k)=J_*R_k-jb_{k-1}(R_{k-1}+R_k)
                  +j(b_k-b_{k-1})m^k,
\]
so \(\|F_h(h^k)\|\le(\|J_*\|+4j)\rho\sqrt n\).
The root estimate of \cite[Lemma~5.7]{SG},
\[
 \|y-y_*\|\le
 \left(1+\frac{\|J_*\|+3j}{b}\right)\|F_h(y)\|
 \quad\hbox{when }\|F_h(y)\|<2\rho_0\sqrt n,
\]
and the Lipschitz property of \(\tanh\) therefore give
\(\|m^k-m_*\|\le C_{\rm lead}\rho\sqrt n\), with a constant
independent of the depth.  For a deterministic unit vector \(e\),
sum \eqref{warm:sqrt-residual} for \(R_1,\ldots,R_k\), with source
\(e\) and multiplier \(\widehat C_k\).  Since
\(\sigma-m^k=\sum_{l\le k}R_l\), this bounds its reweighted mean
against \(e\) by \(C_\rho(1+\sqrt{\mathcal D_h(f)/N_h})\).
The weighted root errors together cost at most
\(C_{\rm lead}\epsilon\sqrt n\), because the cutoffs sum to at most
one.  The remaining deficit costs at most
\(2\sqrt n\,\nu_h(D_{\rm sel}>0)\), controlled by
\eqref{warm:sqrt-deficit}.  Sum over the fixed number of cutoffs,
divide by \(\sqrt n\), and take the supremum over \(e\).
This proves \eqref{warm:sqrt-posterior}.
\end{proof}

\begin{lemma}[Retention at order-\(n\) entropy]\label{warm:large-retention}
Fix \(0<\lambda<C_{\mathrm{lev}}<\infty\), \(C_E<\infty\), and \(0<\omega<1\).
On Gaussian disorder events of probability tending to one, every
test \eqref{warm:restricted-class} with \(\lambda n\le L\le C_{\mathrm{lev}}n\)
and
\begin{equation}
 \frac{\cE(f)}N\le C_E n^{1-\omega}                       \label{warm:large-energy-assumption}
\end{equation}
has an observation path to a product posterior for which
\(H_{\rm end}\ge c_{\lambda,C_{\mathrm{lev}}}n\).  The path and disorder events
are chosen independently of \(f\).
\end{lemma}

\begin{proof}
We use the path and the field tests in
\cite[Section~19.6]{CM}; we verify that its entropy argument remains
valid under \eqref{warm:large-energy-assumption}.  On source disorder
events independent of the initial spin law, the calculation around
Equation~(19.20) there, using its Lemma~17.4, Equation~(17.12), and
Proposition~17.5, gives the following early field-test bound.
On a fixed interval \([t_*,t_b]\) of small
positive scalar precision, the diagnostic and stability failures
have integrated probability \(o(1)\), uniformly under every initial
spin law.  This follows there by conditioning on the spin: the
centered Gaussian-noise diagnostics have integrated failure \(o(1)\),
the additional spin shift has normalized norm \(O(t)\), and the
Gaussian low-band small-ball estimate is uniform in its center.
No factor \(e^L\) is used for this interval.  From \(t_b\) until
the interaction coefficient reaches a fixed globally controlled
range, the buffered stability test fails anywhere on the path with
probability at most \(e^{-a n}\) under the ordinary law \(Q_\mu\),
for a fixed \(a>0\).
This is the later field bound in \cite[Lemma~17.3, Section~19.4, and
Section~19.6 after Equation~(19.21)]{CM}.
The word diagnostics and stability constants on the retained
regions are those of Lemma~\ref{warm:posterior-mean}.

Here is the order of the fixed choices in that construction.  Choose
first a large scalar endpoint \(t_{\rm end}\) and a small interaction
coefficient below which covariance is globally bounded.  Choose
\(t_b>0\) for the small-positive-field estimate, and obtain the
buffer and rate \(a\) for the later test.  The norm bound and the
fixed total duration give constants \(C_{\rm jump}\) and
\(C_{\rm aft}\) for the jump loss divided by \(t_*n\) and the
loss after failure divided by \(n\), uniformly when
\(1-\beta'\le t_*/3\).  Choose a target drift budget \(\eta_1>0\)
so that
\[
 3\eta_1(1+C_{\rm aft}/a)<\lambda/4.
\]
Next choose \(0<t_*\ll t_b\) with
\(C_{\rm jump}t_*\le\eta_1\), and finally \(1-\beta'>0\)
sufficiently small relative to these choices, including
\(1-\beta'\le t_*/3\).  Start with precision
\[
 B_*=t_*I+(1-\beta')W,
\]
which is positive when \(1-\beta'\) is small enough.  Increase scalar
precision to \(t_{\rm end}I+(1-\beta')W\), and then traverse
\[
 B_v=(t_{\rm end}+4v)I+(1-\beta'+v)W,\qquad 0\le v\le\beta'.
\]
The latter rate \(4I+W\) is positive on the norm event.  The endpoint
is \( (t_{\rm end}+4\beta')I+W\), so its posterior is a product.
Before the globally controlled portion, the posterior interaction
coefficient is either \(\beta'<1\) or in a fixed compact subinterval
of \((0,1)\).  Reduce the available diagonal enlargement in
\eqref{warm:stability} to \(0<\eta<(\beta')^{-1}-1\).
Once these stability margins and \(C_{\rm lead}\) are fixed, choose
the mean accuracy \(\epsilon\) in Lemma~\ref{warm:posterior-mean}
so that its leading error, integrated over the bounded path and
converted to a quadratic mean bound, is less than \(\eta_1\).
Then fix its residual tolerance, selection depth, finite word
diagnostics and constant \(C_\epsilon\), and finally take \(n\)
large.

On a valid posterior, apply \eqref{warm:sqrt-posterior} to \(f\)
and to \(1\).  With
\(a_u=\|\nu_u[X]-\mu_u[X]\|/\sqrt n\), this gives
\[
 a_u\le2C_{\rm lead}\epsilon+
       \frac{C_\epsilon}{\sqrt n}
          \left(2+\sqrt{\frac{\cE_{\mu_u}(f)}{N_u}}\right).
\]
By \eqref{warm:conditioned-energy} and Cauchy--Schwarz, the
\(Q_\nu\)-mean of the square root, also with any validity indicator,
is at most \(\sqrt{\cE(f)/N}\).  Under
\eqref{warm:large-energy-assumption} the resulting error is
\[
 \frac{C_\epsilon}{\sqrt n}
       \left(2+\sqrt{\frac{\cE(f)}N}\right)
       =O_\epsilon(n^{-1/2}+n^{-\omega/2})=o(1).
\]
This bound is applied at deterministic observation times; inserting
\(\ind_{\{u<\tau\}}\) for a stopping time only decreases its
Cauchy--Schwarz bound.  Since \(a_u\le2\) always, \(a_u^2\le2a_u\).
The early invalid paths cost at most four times their integrated
probability, hence \(o(1)\).  Thus the normalized quadratic
mean-shift integral, on valid paths and before any later stopping,
is at most \(\eta_1+o(1)\) by the preceding choice of \(\epsilon\).

Let \(u_g\) begin the globally controlled portion and let \(\tau\)
be the first failure of the later buffered test in \([t_b,u_g]\),
with the immediate and no-failure conventions of the source.  The
initial jump leaks at most \(C_{\rm jump}t_*n\) entropy: compare to pure noise
as above and use \(\|B_*\|\le Ct_*\) and \(\|X\|^2=n\).
The stopped form of \eqref{warm:spin-loss}, with the preceding
integrated estimate, gives, with \(Q^\tau\) denoting the law of the
stopped observations, excluding the spin,
\begin{equation}
 D(Q_\nu^\tau\Vert Q_\mu^\tau)
       \le(C_{\rm jump}t_*+\eta_1+o(1))n .              \label{warm:stopped-entropy}
\end{equation}
For the event \(E\) of a failure before \(u_g\), the ordinary
probability is at most \(e^{-an}\).  Data processing for its indicator
then gives
\[
 D(Q_\nu^\tau\Vert Q_\mu^\tau)
 \ge Q_\nu(E)\log\frac1{Q_\mu(E)}-\log2,\qquad
 Q_\nu(E)\le\frac{C_{\rm jump}t_*+\eta_1+o(1)}a .
\]
These are exactly the deductions in
\cite[Equations~(19.21)--(19.22)]{CM}, now with the enlarged energy
range justified by Lemma~\ref{warm:posterior-mean}.
After a failure the remaining precision rate and duration are bounded,
and both spin means have norm at most \(\sqrt n\); the remaining
loss per failed path is therefore at most \(C_{\rm aft}n\).  The total loss
before \(u_g\) is bounded by
\[
 \left(C_{\rm jump}t_*+\eta_1+o(1)\right)
       (1+C_{\rm aft}/a)n .
\]
The fixed choices above make this less than \(\lambda n/4\).
The initial entropy is at least \(L/2\ge\lambda n/2\), hence
\(H(u_g)\ge\lambda n/4\).  In the remaining portion the spin
covariance is bounded for every field.  The linear entropy test
then gives \(-\dot H\le CH\); over its bounded duration it retains
a fixed fraction of \(H(u_g)\).  This proves the lemma.
\end{proof}

\begin{proof}[Proof of Proposition~\ref{warm:profiles}]
Fix \(C_{\mathrm{lev}},\xi\) and a desired disorder confidence.  Apply
Lemma~\ref{warm:small-profile} with
\(\kappa<\min\{\xi,\,5/3-1.662\}\).  After fixing its \(M\), choose
its \(\delta>0\) as in that proof.  For \(L\le n^\delta\),
\eqref{warm:small-profile-bound} and
\eqref{warm:modified-classical} imply both claimed inequalities.

For \(n^\delta<L\le\lambda n\), use the proportional branch of
\cite[Equation~(19.16) and Section~19.4]{CM}.  It gives, at the
product endpoint,
\begin{equation}
 H_{\rm end}\ge cL(L/n)^{(2+2\epsilon)/3}.                \label{warm:moderate-retention}
\end{equation}
Here \(\epsilon>0\) is any fixed sufficiently small tolerance,
and \(\lambda>0\) is chosen afterward.  We justify the exact test
class being used.  The starting dimension in that source is
\(\max\{k_{\min},c_0L\}\), with \(k_{\min}\asymp\sqrt{\log n}\).
It equals \(c_0L\) throughout the present range for large \(n\).
The proportional part of its initial jump estimate is the pure-noise
argument reproduced in Lemma~\ref{warm:small-profile}; it requires only
the entropy at least \(L/2\) and density at most \(e^L\).
The source's support condition is used in its other, minimum-dimension
branch.  Its subsequent propagation uses the likelihood and
conditional-entropy tails in \cite[Equation~(19.8)]{CM} and the
ordinary path tests.  These also require only the same density and
entropy bounds.  The continuation through its finite remaining
scales uses those likelihood tails and conditional Pinsker bounds,
then globally bounded covariance, and so preserves this test class.
Thus \eqref{warm:moderate-retention} applies to
\eqref{warm:restricted-class} without a support restriction.

Since \(L\ge n^\delta\gg\log n\), comparison with the first bound
of \eqref{warm:product-bounds} gives
\[
 \Tn\,\frac{\cE(f)}N
 \ge c n^{-2\epsilon/3}L^{7/6+2\epsilon/3}
 \ge c L^{7/6-(2\epsilon/3)(1/\delta-1)}.
\]
The additive \(O(n^{-3})\) is smaller than
\eqref{warm:moderate-retention}.  The second product bound similarly
gives
\[
 \Tn\,\frac{\cI(f^2)}N
 \ge c n^{-2\epsilon/3}L^{5/3+2\epsilon/3}
 \ge c L^{5/3-(2\epsilon/3)(1/\delta-1)}.
\]
Choose \(\epsilon\) after \(\delta\) so that
\[
 \frac{2\epsilon}{3}(1/\delta-1)
       <\min\{\xi,\,7/6-1.162\}.
\]
Then choose \(\lambda\) below the finitely many fixed rates and
shift buffers in that source argument.  This proves the profiles
in the moderate range.

Finally let \(\lambda n\le L\le C_{\mathrm{lev}}n\).  If the classical bound
were false, then
\(\cE(f)/N\le Cn^{1.162-2/3}\).  This meets
\eqref{warm:large-energy-assumption} for a fixed positive
\(\omega\), whereas the first product bound would give
\[
 H_{\rm end}\le Cn^{1/2+1.162-2/3}+O(n^{-3})=o(n),
\]
contrary to Lemma~\ref{warm:large-retention}.  If the modified bound
were false, \eqref{warm:modified-classical} would instead give
\(\cE(f)/N\le Cn^{1-\xi}\), so that lemma applies with
\(\omega=\xi\).  The second product bound would give
\(H_{\rm end}\le Cn^{1-\xi}=o(n)\), again a contradiction.

All the choices are uniform in \(f\).  More explicitly, the
variance exponent and \(M\) are fixed first, then the small
\(\delta\), then the sharp cap and moderate-entropy tolerances,
then \(\lambda\) and the fixed parameters of the order-\(n\) path.
The finite word length in that last path is fixed after its mean
accuracy.  The final threshold \(K_0\) and \(n\) are chosen last.
Intersecting this finite collection of disorder events preserves
arbitrarily high limiting probability.  This proves the proposition.
\end{proof}

\subsection{Passage from profiles to arbitrary initial laws}
\label{warm:layers-section}

On the norm event \(\|W\|\le3\), comparison of the largest and
smallest spin Hamiltonians gives
\begin{equation}
 \min_\sigma\mu(\sigma)\ge e^{-C_a n}                     \label{warm:atom}
\end{equation}
for an absolute \(C_a\).  In the remainder of the proof choose
\(C_{\mathrm{lev}}>2C_a+1\) in Proposition~\ref{warm:profiles}.  For any density
\(z\ge0\), \(\mu[z]=1\) implies \(z(\sigma)\le1/\mu(\sigma)\);
thus every nonempty positive logarithmic level of \(z\) lies below
\(C_a n\).

We first isolate the bounded test associated with one such level.
This same construction will be used for entropy and higher log moments.

\begin{lemma}[A bounded test from a logarithmic level]\label{warm:log-layer}
Fix \(b>0\) and \(r>0\).  Let \(z>0\) be a density and suppose that,
for some \(d\ge v/2\) and sufficiently large \(v\),
\[
 B=\mu[z\ind_{\{d\le\log z\le1.1d\}}]
       \ge b(v/d)^r>0 .
\]
Define
\begin{equation}
 u=\mathcal U_d(z)
   :=\min\{e^{1.1d},\,z\min(1,ze^{-0.8d})\},
 \qquad f^2=e^{-1.1d}u .                                 \label{warm:level-transform}
\end{equation}
Then \(0\le f\le1\), \(B\le\mu[u]\le1\), and, if
\(N=\mu[f^2]=e^{-L}\),
\[
 1.1d\le L\le1.1d+C_{b,r}\bigl(1+\log_+(d/v)\bigr),
 \qquad
 D\!\left(\frac{f^2\mu}{N}\middle\Vert\mu\right)\ge L/2 .
\]
For large \(v,n\), this test lies in the range
\(K_0\le L\le C_{\mathrm{lev}}n\) of Proposition~\ref{warm:profiles}.  In addition,
\begin{equation}
 \cI(\mathcal U_d(z))\le4\cI(z).                          \label{warm:level-dissipation}
\end{equation}
\end{lemma}

\begin{proof}
In logarithmic coordinates \(l=\log z\), the logarithm of
\(\mathcal U_d(e^l)\) is
\[
 2l-0.8d\quad(l\le0.8d),\qquad
 l\quad(0.8d\le l\le1.1d),\qquad
 1.1d\quad(l\ge1.1d).
\]
Consequently \(u\le z\), and \(u=z\) on the indicated level.
This proves \(B\le\mu[u]\le1\) and
\(L=1.1d-\log\mu[u]\), with the stated upper bound from the lower
bound on \(B\).

Let \(E=\{\log z\ge0.7d\}\).  Since \(u\le e^{-0.1d}z\) on its
complement,
\[
 \frac{\mu[u\ind_{E^c}]}{\mu[u]}\le\frac{e^{-0.1d}}B,
 \qquad \mu(E)\le e^{-0.7d}.
\]
Set \(e_d=e^{-0.1d}/B\).  Data processing for the indicator of
\(E\) gives
\[
 D\!\left(\frac{u\mu}{\mu[u]}\middle\Vert\mu\right)
 \ge (1-e_d)\log\frac1{\mu(E)}-\log2
 \ge0.7d(1-e_d)-\log2 .
\]
Uniformly for \(d\ge v/2\), the hypothesis on \(B\) makes
\(e_d\to0\) and \(1+\log_+(d/v)=o(d)\) as \(v\to\infty\).
The last lower bound therefore exceeds \(L/2\) for large \(v\).
The level is nonempty, so \(d\le C_an\) by \eqref{warm:atom};
the bound on \(L\), with \(C_{\mathrm{lev}}>2C_a+1\), puts it below \(C_{\mathrm{lev}}n\)
for large \(n\).  Its lower bound puts it above \(K_0\) for large
\(v\).

Finally \(\mathcal U_d\) is nondecreasing and \(2\)-Lipschitz as a
function of \(z\), while \(\log\mathcal U_d(e^l)\) is nondecreasing
and \(2\)-Lipschitz as a function of \(l\).  The product of their
increments on an edge is at most four times the product of the
increments of \(z\) and \(\log z\).  Summing proves
\eqref{warm:level-dissipation}.
\end{proof}

\begin{lemma}[Entropy dissipation for every density]\label{warm:entropy-profile}
With tight constants, every strictly positive density \(z\) with
\(\Ent_\mu(z)\ge K_1\) satisfies
\begin{equation}
 \Tn\cI(z)\ge c\,\Ent_\mu(z)^{1.662}.                    \label{warm:full-entropy-profile}
\end{equation}
\end{lemma}

\begin{proof}
Put \(H=\Ent_\mu(z)\) and \(d_j=(H/2)(1.1)^j\).
The part of \(\mu[z\log z]\) with \(\log z<H/2\) is at most \(H/2\).
Thus the levels \([d_j,1.1d_j]\) account for at least \(H/2\).
If their tilted masses \(B_j\) were all less than
\(c(H/d_j)^{1.1}\), their contribution would be at most
\[
 1.1\sum_{j\ge0}d_jB_j
 <1.1cH^{1.1}\sum_{j\ge0}d_j^{-0.1}\le CcH .
\]
Choose \(c\) small.  One level therefore has
\(B\ge c(H/d)^{1.1}\).  Apply Lemma~\ref{warm:log-layer} with \(v=H\)
and \(r=1.1\), followed by \eqref{warm:modified-profile} with any
fixed \(\xi<5/3-1.662\).  Homogeneity of \(\cI\) gives
\[
 \Tn\cI(z)\ge\tfrac14\Tn\cI(u)
       \ge c\,\mu[u]L^{1.662}
       \ge cB d^{1.662}
       \ge c'H^{1.1}d^{0.562}\ge c''H^{1.662}.
\]
The constants and the required lower threshold on \(H\) are uniform.
\end{proof}

For the kernel we need a classical support profile.  The next
deduction is the equal-height and dyadic-layer argument of
\cite[Lemma~19.1]{CM}, keeping the two powers from
\eqref{warm:classical-profile}.

\begin{lemma}[A support profile]\label{warm:support-profile}
With the constants and \(\delta\) of Proposition~\ref{warm:profiles}, every
nonnegative \(g\), supported on a set of \(\mu\)-mass \(e^{-L_{\rm supp}}\)
with \(L_{\rm supp}\ge K_0\), satisfies
\begin{equation}
 \Tn\cE(g)\ge c\mu[g^2]\,\Phi_n(L_{\rm supp}),\qquad
 \Phi_n(L)=
 \begin{cases}
 \min\{L^{1.654},\,n^{1.154\delta}\},&L\le n^\delta,\\
 L^{1.154},&L>n^\delta .
 \end{cases}                                                \label{warm:support-bound}
\end{equation}
The zero function satisfies the assertion by convention.
\end{lemma}

\begin{proof}
First let \(0\le g\le1\) equal one on a nonempty set of mass
\(a=e^{-L_a}\), and let its support have mass at most \(e^{-L_{\rm supp}}\).
Take
\[
 K=\left\lceil\log_2(1+L_a/L_{\rm supp})\right\rceil,\qquad
 g_i=K\bigl((g-i/K)_+\wedge K^{-1}\bigr),\quad 0\le i<K .
\]
Let \(S_i=\mu(g>i/K)\), \(T_i=\mu(g\ge i/K)\).
Some \(i\) satisfies \(T_{i+1}\ge S_i^2\); otherwise induction gives
\(a\le T_K<e^{-2^KL_{\rm supp}}<a\).  For this \(i\), with \(N_i=\mu[g_i^2]\),
\[
 a\le T_{i+1}\le N_i\le S_i,\qquad
 D\!\left(\frac{g_i^2\mu}{N_i}\middle\Vert\mu\right)
       \ge\log(1/S_i)\ge\tfrac12\log(1/N_i).
\]
Its logarithmic inverse mass \(L_i\) lies in \([L_{\rm supp},L_a]\),
and \(L_a\le C_an\) by \eqref{warm:atom}.  For \(L\ge K_0\)
large enough, the function
\[
 e^{-L}
 \begin{cases}L^{1.662},&L\le n^\delta,\\
               L^{1.162},&L>n^\delta
 \end{cases}
\]
is decreasing within each range and has a downward jump at their
boundary.  Since \(\cE(g_i)\le K^2\cE(g)\),
\eqref{warm:classical-profile} implies
\[
 \Tn\cE(g)\ge
 \frac{ca}{(1+\log L_a)^2}
 \begin{cases}L_a^{1.662},&L_a\le n^\delta,\\
               L_a^{1.162},&L_a>n^\delta .
 \end{cases}
\]
Absorbing the logarithmic denominator by lowering each exponent by
\(0.008\) gives
\begin{equation}
 \Tn\cE(g)\ge ca
 \begin{cases}L_a^{1.654},&L_a\le n^\delta,\\
               L_a^{1.154},&L_a>n^\delta .
 \end{cases}                                                \label{warm:set-capacity}
\end{equation}
Indeed \(L^{0.008}/(1+\log L)^2\) has a positive infimum for
\(L\ge K_0\).

For arbitrary nonnegative \(g\) with the stated support, use
\(\varphi_j=(g-2^j)_+\wedge2^j\), \(j\in\mathbb Z\), and let
\(a_{j+1}=\mu(g\ge2^{j+1})\).  Omit empty upper level sets.
The normalized layer \(2^{-j}\varphi_j\) is one on that set and
has the original support restriction.  Its \(L_a=\log(1/a_{j+1})\)
is at least \(L_{\rm supp}\).  The right side of
\eqref{warm:set-capacity}, divided by \(a\), is at least
\(\Phi_n(L_{\rm supp})\): before the boundary it is at least \(L_{\rm supp}^{1.654}\),
and after the boundary it is at least \(n^{1.154\delta}\), or at
least \(L_{\rm supp}^{1.154}\) when \(L_{\rm supp}>n^\delta\).
For each oriented edge the nonnegative increments of the
\(\varphi_j\)'s sum to the increment of \(g\), so the sum of their
squares is at most its square.  Therefore
\[
 \Tn\cE(g)\ge\sum_j\Tn\cE(\varphi_j)
       \ge c\Phi_n(L_{\rm supp})\sum_j4^ja_{j+1}.
\]
Pointwise,
\(\sum_j4^j\ind_{\{g\ge2^{j+1}\}}\ge g^2/12\).
Integration proves \eqref{warm:support-bound}.
\end{proof}

\begin{proof}[Proof of Theorem~\ref{thm:warming}]
Fix the parameters in the theorem and intersect the finitely many
events needed for the preceding results.  The finite heat-bath chain
has a strictly positive kernel at every positive time.  Hence
\(w_s^\rho>0\) for \(s>0\), for every initial law, and symmetry of
the generator gives
\[
 \frac{\dd}{\dd s}\Ent_\mu(w_s^\rho)
       =-\Tn\cI(w_s^\rho).
\]
When the entropy exceeds the threshold of
Lemma~\ref{warm:entropy-profile}, its derivative is at most
\(-c\Ent_\mu(w_s^\rho)^{1.662}\).  Integrating this differential
inequality from any earlier positive time, and then letting that
time decrease to zero, yields uniformly in the initial law
\[
 \Ent_\mu(w_s^\rho)
       \le C\left(1+s^{-1/0.662}\right)
       \le C'\left(1+s^{-1.54}\right).
\]
This proves \eqref{warm:entropy}.

For the kernel, put \(Q_s=\mu[(w_s^\rho)^2]=e^{L_s}\).
The atom bound gives \(Q_s\le e^{C_an}\), so its logarithm is
within the range where the support profile was proved.
Whenever \(L_s\) is large, set
\(g=(w_s^\rho-Q_s/4)_+\).  Its support has mass at most
\(4e^{-L_s}\) by \(\mu[w_s^\rho]=1\), and
\[
 \mu[g^2]\ge Q_s-2(Q_s/4)\mu[w_s^\rho]=Q_s/2,\qquad
 \cE(g)\le\cE(w_s^\rho).
\]
Its actual support logarithm is
\[
 L_{\rm supp}=\log\frac1{\mu\{w_s^\rho>Q_s/4\}}
       \ge L_s-\log4.
\]
The support is nonempty by the preceding square-mass bound, so
\eqref{warm:atom} also gives \(L_{\rm supp}\le C_an\).
Since \(\dot Q_s=-2\Tn\cE(w_s^\rho)\),
Lemma~\ref{warm:support-profile} and monotonicity of \(\Phi_n\) give
\[
 -\dot L_s\ge c\,\Phi_n(L_s-\log4)
\]
above a fixed threshold.  A shift of the argument by \(\log4\)
only changes constants there.  For any large \(K'\), the time
allowed by this differential inequality above \(K'\) is at most
\begin{equation}
 \int_{K'}^\infty\frac{C\,\dd L}{\Phi_n(L)}
       \le C(K')^{-0.654}+C n^{-0.154\delta}.
                                                               \label{warm:hitting-integral}
\end{equation}
For completeness, below
\(n^{\delta(1.154/1.654)}\) integrate \(L^{-1.654}\).
The following constant portion of \(\Phi_n\), up to \(n^\delta\),
costs at most \(n^\delta/n^{1.154\delta}\), and the integral of
\(L^{-1.154}\) above \(n^\delta\) has the same order.
Inverting \eqref{warm:hitting-integral} gives
\[
 L_s\le C(1+s^{-1/0.654})
       \qquad\hbox{for }s\ge Cn^{-0.154\delta},
\]
uniformly even for point initial laws.  Reversibility and two half
steps show that
\[
 p_{s\Tn}(\sigma,\tau)
 =\mu[p_{s\Tn/2}(\sigma,\cdot)p_{s\Tn/2}(\tau,\cdot)]
 \le\bigl(\mu[p_{s\Tn/2}(\sigma,\cdot)^2]\,
          \mu[p_{s\Tn/2}(\tau,\cdot)^2]\bigr)^{1/2}.
\]
The preceding bound, with \(1/0.654<1.54\), proves
\eqref{warm:kernel} after enlarging the vanishing threshold.

We finish by proving \eqref{warm:log-moment}.  For the fixed integer
\(m\ge1\), let \(\Psi_m\) be the convex function on \((0,\infty)\)
determined by
\[
 \Psi_m(1)=\Psi_m'(1)=0,\qquad
 \Psi_m''(z)=\frac{(1+|\log z|)^{m-1}}{z}.
\]
It extends continuously to zero: in its twice-integrated formula the
possible singularity is bounded by
\(\int_0^1(1+|\log z|)^{m-1}\dd z<\infty\).
For a density \(z\), put
\[
 V_m(z)=\mu[\Psi_m(z)],\qquad
 \mathsf M_m(z)=\left(1+\mu[z(\log_+z)^m]\right)^{1/m}.
\]
There are constants depending only on \(m\) such that
\begin{equation}
 c_m\mathsf M_m(z)^m-C_m
       \le V_m(z)\le C_m\mathsf M_m(z)^m.                 \label{warm:moment-comparison}
\end{equation}
Indeed for \(z\ge1\),
\(\Psi_m'(z)=((1+\log z)^m-1)/m\).  Integrating this on
\([1,z]\) gives upper and lower bounds of order
\(z(\log z)^m\) when \(\log z\) is large; the remaining range
contributes at most a constant times \(1+z\).  On \([0,1]\)
\(\Psi_m\) is bounded.

For a strictly positive density \(z\), define the corresponding
dissipation \(\mathcal D_m(z)=\cE(z,\Psi_m'(z))\).
We claim that, for \(\mathsf M_m(z)\) above a fixed threshold,
\begin{equation}
 \Tn\mathcal D_m(z)
       \ge c_{m,\xi}\mathsf M_m(z)^{m+2/3-\xi}.
                                                               \label{warm:moment-dissipation}
\end{equation}
Put \(V=\mathsf M_m(z)\) and use levels
\(d_j=(V/2)(1.1)^j\).  If \(V\) is large, the part of
\(\mu[z(\log_+z)^m]\) below \(V/2\) is at most \(2^{-m}V^m\).
Geometric summation as in Lemma~\ref{warm:entropy-profile} gives a level
\(d\ge V/2\) with
\[
 B=\mu[z\ind_{\{d\le\log z\le1.1d\}}]
       \ge c_m(V/d)^{m+0.1}.
\]
Indeed otherwise
\((1.1)^m\sum_jd_j^mB_j\le C_m c_m V^m\) would be too small.
The test \(u=\mathcal U_d(z)\) from Lemma~\ref{warm:log-layer} therefore
satisfies \eqref{warm:restricted-class}, and
\eqref{warm:modified-profile} gives
\begin{equation}
 \Tn\cI(u)\ge c\,\mu[u]L^{5/3-\xi}
       \ge cB d^{5/3-\xi}.                               \label{warm:moment-layer-profile}
\end{equation}

We check that the new dissipation controls this layer with its
additional weight \(d^{m-1}\).  On an edge write
\(a\ge b>0\), \(x=\log a\), \(y=\log b\).  Then
\[
 \Psi_m'(a)-\Psi_m'(b)
       =\int_y^x(1+|t|)^{m-1}\dd t.
\]
If \(x\ge0.4d\), the last integral is at least
\(c_m d^{m-1}(x-y)\).  To verify this uniformly in \(y\), use the
top half of \([0,x]\) when \(-x\le y\le x/2\), the entire interval
when \(y\ge x/2\), and the negative half nearest \(y\) when
\(y<-x\); in each case a fixed fraction of the interval has
\(|t|\) at least a fixed fraction of \(x\).  The two
Lipschitz bounds in the proof of Lemma~\ref{warm:log-layer} then apply.
If \(x<0.4d\), both endpoints are in the lower branch of
\(\mathcal U_d\).  On \([b,a]\) its derivative is at most
\(2e^{-0.4d}\), while its logarithm has slope two in log
coordinates.  Since \(d^{m-1}e^{-0.4d}\le C_m\), the same
conclusion follows using \((1+|t|)^{m-1}\ge1\).
Thus in both cases
\[
 d^{m-1}\bigl(\mathcal U_d(a)-\mathcal U_d(b)\bigr)
       \bigl(\log\mathcal U_d(a)-\log\mathcal U_d(b)\bigr)
 \le C_m(a-b)\bigl(\Psi_m'(a)-\Psi_m'(b)\bigr).
\]
Summation gives
\(\mathcal D_m(z)\ge c_m d^{m-1}\cI(u)\).
Combining this with \eqref{warm:moment-layer-profile} and the
lower bound on \(B\), and using \(2/3-\xi>0.1\), gives
\[
 \Tn\mathcal D_m(z)
 \ge cB d^{m+2/3-\xi}
 \ge c_m V^{m+0.1}d^{2/3-\xi-0.1}
 \ge c_m' V^{m+2/3-\xi}.
\]
This proves \eqref{warm:moment-dissipation}.

Along the finite semigroup at positive times,
\[
 \frac{\dd}{\dd s}V_m(w_s^\rho)
       =-\Tn\mathcal D_m(w_s^\rho).
\]
Equations \eqref{warm:moment-comparison}--\eqref{warm:moment-dissipation}
give, above a fixed threshold,
\[
 \frac{\dd}{\dd s}V_m(w_s^\rho)
       \le-c V_m(w_s^\rho)^{1+(2/3-\xi)/m}.
\]
Integrating from an earlier positive time and letting that time
decrease to zero yields
\(V_m(w_s^\rho)\le C(1+s^{-m/(2/3-\xi)})\).
The comparison \eqref{warm:moment-comparison} proves
\eqref{warm:log-moment}.  This argument allows initial zeros,
because \(\Psi_m\) is continuous at zero and the kernel is strictly
positive afterward.

The three conclusions \eqref{warm:entropy}--\eqref{warm:log-moment}
depend only on the off-diagonal couplings.  Let \(\mathcal G_n^{\rm aug}\)
be the final event in the augmented Gaussian space, and put
\[
 \mathcal G_n^{\rm off}
 =\left\{W_{\rm off}:
       \Prob(\mathcal G_n^{\rm aug}\mid W_{\rm off})>0\right\}.
\]
Then \(\Prob(\mathcal G_n^{\rm off})\ge
\Prob(\mathcal G_n^{\rm aug})\).  For fixed off-diagonal couplings the
three inequalities are deterministic, so their validity on an augmented
realization in \(\mathcal G_n^{\rm aug}\) descends to
\(\mathcal G_n^{\rm off}\).  Taking this as the
theorem event proves the stated disorder quantifiers.
\end{proof}

\begin{corollary}[A polynomial tail bound for the heat-kernel logarithm]
\label{warm:kernel-log-tail}
Fix \(\beta>0\), \(A>0\), \(0<C_\star<\infty\), and
\(\vartheta>0\).  On Gaussian disorder events of limiting lower
probability at least \(1-\vartheta\), with a fixed constant \(C\),
\begin{equation}
 \sup_{\sigma}\ \sup_{1\le d\le C_\star\Tn}
 \sum_\tau P_d(\sigma,\tau)
 \ind_{\left\{\log p_d(\sigma,\tau)>
       \frac12 n^\beta(\Tn/d)^{3/2}\right\}}
 \le Cn^{-A}                                                \label{warm:kernel-tail}
\end{equation}
for all sufficiently large \(n\).  For any initial law \(\rho\),
averaging the displayed row probabilities over \(\sigma\sim\rho\)
preserves the same bound.
\end{corollary}

\begin{proof}
Choose \(\xi>0\) so small that
\[
 \frac{\xi}{2/3-\xi}<\frac{\beta}{2},
\]
then choose a fixed integer \(m>2A/\beta\).  Apply
\eqref{warm:log-moment} to the point initial law at \(\sigma\)
and time \(s=d/\Tn\).  Set \(v=\Tn/d\).  Uniformly for
\(1\le d\le C_\star\Tn\), so \(C_\star^{-1}\le v\le n^{2/3}\),
\[
 \frac{1+v^{1/(2/3-\xi)}}{n^\beta v^{3/2}}
 \le C_{\xi,C_\star}
       n^{-\beta+\xi/(2/3-\xi)}
 \le C_{\xi,C_\star}n^{-\beta/2}.
\]
Markov's inequality under the transition law
\(P_d(\sigma,\tau)=\mu(\tau)p_d(\sigma,\tau)\) now bounds the
left side of \eqref{warm:kernel-tail} by
\[
 \left(
 \frac{2\{1+\mu[p_d(\sigma,\cdot)
                  (\log_+p_d(\sigma,\cdot))^m]\}^{1/m}}
      {n^\beta(\Tn/d)^{3/2}}\right)^m
 \le C_m n^{-m\beta/2}\le C_m n^{-A}.
\]
The choices of \(\xi,m\) are fixed in that order before \(n\).
The statement for \(\sigma\sim\rho\) follows by averaging these
row-specific probabilities.
\end{proof}
